\section{Method details}
\label{S-sec:s-method}

The method in \main{\cref*{sec:method}} takes $N$ putative matches $(p_i,q_i)\in\mathbb{R}^3\times\mathbb{R}^3$ and two observed point clouds as input, and returns a rigid motion $(R,t)$. The \ngr{} frontend generates pose hypotheses from the matches. Verification refines several candidates and selects one pose; it can accept hypotheses from any frontend.

\paragraph{Seeded nucleation.} The seeds are geometric motifs of the zero-diagonal compatibility graph $\Lambda$ of \main{\cref*{eq:compat}}. The kernel width is $\sigma_c=0.375\tau$ (3.75\,cm indoors). Pairs are ranked by the second-order compatibility $\Lambda_{ab}(\Lambda^2)_{ab}$ of SC$^2$-PCR~\cite{chen2022sc2} and triples by $T_{abc}=\Lambda_{ab}\Lambda_{bc}\Lambda_{ca}\sum_d\Lambda_{ad}\Lambda_{bd}\Lambda_{cd}$, with a non-degeneracy test on both triangles, $(p_a,p_b,p_c)$ and $(q_a,q_b,q_c)$. Each selected motif seeds one phase, a probability vector over all matches: $w_i^{(0)}\propto\Lambda_{ai}\Lambda_{bi}+\delta_{ia}+\delta_{ib}$ for a pair $(a,b)$ and $w_i^{(0)}\propto\Lambda_{ai}\Lambda_{bi}\Lambda_{ci}+\delta_{ia}+\delta_{ib}+\delta_{ic}$ for a triple $(a,b,c)$. Every match therefore starts with positive weight. A background-load scheduler computes $(N-2)\bar\Lambda^2$, the number of supporting triangles a pair expects from the background ($\bar\Lambda$ is the mean off-diagonal compatibility). Below one it uses 256 pair seeds and the count readout. Otherwise a shortlist of 4096 pairs is extended to $K=1024$ triples with the MAE readout. All $K$ phases evolve independently.

\paragraph{Candidate crowding.} Ranking motifs by compatibility alone allows a single cluster to occupy the entire seed budget. A cluster of $n$ mutually compatible wrong matches has $\binom{n}{2}$ internal pairs, each supported by about $n-2$ triangles (4095 pairs with support near 89 for $n=91$), whereas among $m$ true inliers a true pair collects about $m-2$, only 8--23 for the 10--25 inliers typical of hard pairs. A few such clusters can then take most of the pair shortlist and the nuclei, and the pool concentrates on a few poses (\main{\cref*{sec:crowding}}; \cref{S-sec:s-crowding}).

\paragraph{Finite site capacity.} In competitive adsorption a surface site holds a limited amount of adsorbate, so a dominant species saturates its own sites and leaves the rest for others~\cite{langmuir1918}. Every match $x$ carries the site vacancy $\theta_x=1/(1+n_x/q)$ of \main{\cref*{eq:vacancy}}, where $n_x$ counts the already selected motifs that contain $x$. We discount a candidate motif's affinity by the vacancies of its matches, $T_{abc}\,\theta_a\theta_b\theta_c$ (pairs analogously). The first half of the pair shortlist and of the nuclei keep the original ranking. The second half is filled greedily by discounted affinity from a deeper ranking (65\,536 pairs, 16\,384 triples), counting every selected motif as occupying its sites. The capacities $q=8$ for pairs and $q=4$ for triples were fixed when the variant was pre-registered and never tuned (\cref{S-sec:s-protocol}). Growth, fitting and scoring are unchanged; the resulting frontend is \ls{}. The vacancy is an occupancy feedback in Langmuir form, not a derived adsorption equilibrium.

\paragraph{Relation to diversity selection.} The vacancy penalizes a motif according to how often its matches have already been selected. The penalty is cumulative and soft. Other diversity rules act differently on the same ranked list. Maximal marginal relevance~\cite{carbonell1998mmr} on the match sets, with its default trade-off $\lambda=\frac12$, subtracts the largest overlap with any single selected motif. A hard capacity rejects a motif once one of its matches is held $q$ times (a match may sit in at most 8 pairs or 4 triples); when no admissible motif is left, the remaining nuclei are taken in ranking order. Non-maximum suppression in pose space rejects a nucleus whose three-point pose repeats an earlier one. \main{\cref*{sec:crowding}} and \cref{S-tab:s-diversity} compare the four rules with the same budget, first halves and deep rankings. The two soft rules on match reuse tie; both reach a higher ceiling than hard capacity and pose suppression.

\paragraph{Total-reflux growth, refits and pool score.} Each phase is grown for ten stages by \main{\cref*{eq:growth}} with $\epsilon=1$ and $\beta=2$. The factor $w_i^{(s)}$ multiplies the support by the current composition and retains the enrichment from the previous stage. All matches stay available; only the number of phases is capped. Each grown phase gives a weighted Kabsch fit~\cite{kabsch1976}, followed by two hard refits with residual radius $\tau_c=\tau$ (0.04 in the frontend's internal units, where $\tau$ maps to 0.04). The triple branch scores the pool by $\sum_i[1-r_i/\tau_c]_+$ on the matches, where $r_i$ is the residual of match $i$ under the hypothesis.

\paragraph{Multi-stream finishing.} Finishing keeps the frontend's native winner, shortlists geometrically distinct candidates ($5^\circ$ and $\tau$ pose suppression), and fills at most $B=32$ refinements: up to 16 primary modes allocated by interleaved rankings of three supports, then extra modes from the non-dominated set. Each selected candidate is refined by point-to-plane ICP~\cite{besl1992icp,chen1992medioni} on the observed surfaces at radii $2\tau$ and $\tau$ ($\tau=0.1$\,m indoors, 0.3\,m on ETH, 0.6\,m on KITTI). For each refined pose $h$ it records the truncated correspondence support $M(h)$, the bidirectional surface support $G(h)$, and an information-weighted support $I(h)$ whose query weights are the leverages $\ell_i=J_i(\sum_jJ_j^\top J_j)^{\dagger}J_i^\top$ of the point-to-plane rows $J_i=[(\widetilde x_i\times n_i)^\top,n_i^\top]$~\cite{gelfand2003stable}, with $\widetilde x_i$ the centered and scale-normalized query point and $n_i$ its normal. The weights are computed once on each cloud's query points and do not depend on the pose. The finishing-only output keeps the candidate with the largest $M\sqrt{G}$ as an anchor. A challenger replaces it only if it raises $M\sqrt{I}$ without lowering $I$. This anchor rule is the baseline over which the release assay had to pass its pre-registered margin (\cref{S-sec:s-protocol}). The same finishing procedure runs after every frontend, without labels or learned weights.

\paragraph{Release assay: two specifications.} The assay checks two specifications of each refined candidate $k$ in the refined pool $V$. The \emph{support} channel measures how much of the scene the pose explains, $s_k=M_k\sqrt{I_k}/\max_{j\in V}M_j\sqrt{I_j}$. The \emph{shape} channel measures how tightly the supported points lie on the other surface. For the $m$ query points whose nearest-neighbor distance $d$ is below $\tau$, let $u=d^2/\tau^2$ and $\bar u=\frac1m\sum u$.

\paragraph{The null and the Cram\'er rate.} We choose the idealized radial reference $p_0(d)=2d/\tau^2$ on $[0,\tau)$, motivated by uniform area in a planar disc. It gives $\Pr(d\le r\mid d<\tau)=r^2/\tau^2$. Hence $u$ is uniform on $[0,1)$, $\bar u\approx\frac12$, and $\mathbb{E}\,e^{-\beta u}=(1-e^{-\beta})/\beta$. The Cram\'er rate of a smaller mean under this null is the Legendre transform of $\log\mathbb{E}\,e^{-\beta u}$,
\begin{equation}
\psi(\bar u)=\sup_{\beta\ge0}\Big\{-\beta\bar u-\log\frac{1-e^{-\beta}}{\beta}\Big\},
\label{S-eq:s-psi}
\end{equation}
which is \main{\cref*{eq:psi}}. If the $m$ null residuals are independent, the probability that their mean is at most $\bar u<\frac12$ is bounded by $e^{-m\psi(\bar u)}$ (Chernoff's bound). As $m$ grows, this probability decays at that rate (Cram\'er's theorem), so $\psi$ measures per-point evidence of residual concentration. The objective in braces is concave in $\beta$, vanishes at $\beta=0$ and has slope $\frac12-\bar u$ there, so $\psi$ is zero for $\bar u\ge\frac12$ and positive below. Averaging $\psi$ over the two directions (source queries against the target surface and target queries against the source surface) gives $J_k$, and $h_k=\exp(J_k-\max_{j\in V}J_j)$. The assay releases $\hat k=\arg\max_{k\in V}\min\{s_k,h_k\}$ (\main{\cref*{eq:assay}}), breaking ties by $s_k$, then $h_k$, then queue order. The definitions fix the constants without fitting.

\paragraph{Minimax-regret interpretation.} Consider the surrogate utility $U_\lambda(k)=\lambda s_k+(1-\lambda)h_k$ of candidate $k$ for an unknown exchange rate $\lambda\in[0,1]$. The regret of choosing $k$ is
\begin{equation}
r_k(\lambda)=\max_{j\in V}U_\lambda(j)-U_\lambda(k).
\label{S-eq:s-regret}
\end{equation}
Its first term is a maximum of functions affine in $\lambda$ and therefore convex, and its second term is affine, so $r_k$ is convex on $[0,1]$ and attains its maximum at an endpoint. Both channels are normalized to a maximum of one over $V$ ($\max_j s_j=1$ by definition and $\max_j h_j=e^0=1$), so $r_k(0)=1-h_k$ and $r_k(1)=1-s_k$. The worst-case regret is the larger of the two, $\max\{1-s_k,1-h_k\}$:
\begin{equation}
\max_{\lambda\in[0,1]}r_k(\lambda)=1-\min\{s_k,h_k\}.
\label{S-eq:s-worst}
\end{equation}
Minimizing the worst-case regret over $k\in V$ is therefore exactly the release rule of \main{\cref*{eq:assay}}. The guarantee concerns this family of surrogate utilities, not registration success itself, and the null behind $\psi$ is an idealization. On real scans, the null is not calibrated, although $\psi$ helps rank candidates. The product and harmonic mean of the two channels perform similarly to the release rule. A weighted sum with one weight fitted on the training cohort succeeds on about 0.5\% more frontend-pair evaluations (\main{\cref*{sec:verification}}; \cref{S-sec:s-verification}).

\paragraph{Optional learned assay.} The learned assay is a conditional logit over the candidates of a pool (a softmax over the pool of a linear score of each candidate's features) with 26 label-free features of the finishing quantities, trained only on the 3DMatch training scenes. We chose its features, regularization and weights on the training cohort. The candidate features had been listed before the protocol was in place (\cref{S-sec:s-protocol}). We report this variant separately from all training-free results.
